\documentclass[11pt,a4paper]{article}
% Shared preamble for the English probability and statistics documents.
% Based on the original course preamble. Compile with pdfLaTeX.
\usepackage[top=1.5cm,bottom=2cm,left=2cm,right=2cm]{geometry}
\usepackage[T1]{fontenc}
\usepackage{lmodern}
\usepackage[utf8]{inputenc}
\usepackage[english]{babel}

\usepackage{amsmath}
\usepackage{amssymb}
\usepackage{amsthm}
\usepackage{mathrsfs}
\usepackage{mathtools}
\usepackage{etoolbox}
\usepackage{tikz}
\usetikzlibrary{shapes.multipart,arrows.meta}
\usepackage{xcolor}
\usepackage{enumerate}
\usepackage[shortlabels]{enumitem}
\usepackage{hyperref}
\usepackage{version}
\usepackage{bbold}
\usepackage{fontawesome5}
\usepackage{fancyhdr}
\lfoot{}
\usepackage{natbib}
\usepackage{graphicx}
\usepackage{pdfpages}
\usepackage{authblk}
    
\usepackage{blindtext}
\usepackage{multicol}
\usepackage{collcell}
\usepackage{wrapfig}
\usepackage{blkarray}
\usepackage{cancel}
\usepackage{appendix}
\usepackage{float}
\makeatletter
\input{tex/assets/slashbox.sty}
\makeatother
\hypersetup{unicode,colorlinks,citecolor=red,filecolor=black,
  linkcolor=red,urlcolor=cyan}
\setlength{\columnsep}{1cm}
\pagenumbering{gobble}

\def\numtext#1{%
  \ifcase#1\or one\or two\or three\or four\or five\or six\or seven\or eight\or nine\or ten\or
  eleven\or twelve\or thirteen\or fourteen\or fifteen\or sixteen\or seventeen\or eighteen\or nineteen\or twenty\or Lots\fi}

\newcommand{\Byte}[1]{%
  \def\nodecontents{}%
  \foreach[count=\l] \k in {#1} {\xappto{\nodecontents}{\noexpand\nodepart{\numtext{\l}}\k}}%
  \node[rectangle split,rectangle split parts=8,rectangle split horizontal,draw] at (2,2){\nodecontents};}

\newcommand{\D}{{\mathrm{d}}}
\newcommand{\cE}{{\mathcal{E}}}
\newcommand{\cG}{{\mathcal{G}}}
\newcommand{\cX}{{\mathcal{X}}}
\newcommand{\cY}{{\mathcal{Y}}}
\newcommand{\cB}{{\mathcal{B}}}
\newcommand{\cF}{{\mathcal{F}}}
\newcommand{\cO}{{\mathcal{O}}}
\newcommand{\cS}{{\mathcal{S}}}
\newcommand{\cT}{{\mathcal{T}}}
\newcommand{\given}{\, \big| \,}
\newcommand{\cL}{{\mathscr{L}}}
\newcommand{\cD}{{\mathcal{D}}}
\newcommand{\cA}{{\mathscr{A}}}
\newcommand{\cP}{{\mathscr{P}}}
\newcommand{\cR}{{\mathscr{R}}}
\newcommand{\A}{{\mathcal A}}
\newcommand{\one}{\mathbb{1}}
\newcommand{\mathbbm}[1]{\mathbb{#1}}

\providecommand{\C}{}
\renewcommand{\C}{\mathbb C}
\newcommand{\R}{\mathbb R}
\newcommand{\Q}{\mathbb Q}
\newcommand{\N}{\mathbb N}
\newcommand{\T}{\mathbb T}
\newcommand{\F}{\mathbb F}
\renewcommand{\P}{\mathbb P}
\newcommand{\Z}{\mathbb Z}
\newcommand{\E}{\mathbb {E}}
\newcommand{\Pbb}[1]{\P\left( #1\right)}
\newcommand{\Ebb}[1]{\E\left[ #1\right]}

\newcommand{\lb}{\left (}
\newcommand{\rb}{\right )}
\newcommand{\lbb}{\left [}
\newcommand{\rbb}{\right ]}
\newcommand{\labs}{\left |}
\newcommand{\rabs}{\right |}
\newcommand{\lbrb}[1]{\lb #1 \rb}
\newcommand{\lbbrbb}[1]{\lbb#1\rbb}
\newcommand{\labsrabs}[1]{\labs#1\rabs}
\newcommand{\lbbrb}[1]{\lbb#1\rb}
\newcommand{\lbrbb}[1]{\lb#1\rbb}
\newcommand{\langlerangle}[1]{\langle#1\rangle}
\newcommand{\lbcurly}{\{}
\newcommand{\rbcurly}{\}}
\newcommand{\curly}[1]{\lbcurly#1\rbcurly}
\newcommand{\abs}{\labsrabs}
\newcommand{\ind}[1]{1_{#1}}
\newcommand{\indep}{\perp \!\!\! \perp}
\renewcommand{\d}{\, \mathrm{d}}

\newtheorem{theorem}{Theorem}
\newtheorem{problemCnt}{Lemma}
\theoremstyle{definition}
\newtheorem{problem}[problemCnt]{Problem}
\newtheorem{problembook}[problemCnt]{\faBook \,  Problem}
\newtheorem{problempencil}[problemCnt]{\faPencil* \,  Problem}
\newtheorem{problemstar}[problemCnt]{\faStar\,  Problem}
\newtheorem{problemstarpencil}[problemCnt]{\faStar \, \faPencil* \,   Problem}
\newtheorem{probleminteresting}[problemCnt]{\faSmile[regular] \,   Problem}
\newtheorem{solution}{Solution}

\newcommand{\fwcell}[1]{\makebox[\arraystretch\normalbaselineskip]{$#1$}}
\newenvironment{enumeratei}{\begin{enumerate}[(i)]}{\end{enumerate}}
\newenvironment{enumeratea}{\begin{enumerate}[(a)]}{\end{enumerate}}
\newenvironment{enumeraten}{\begin{enumerate}[\upshape 1.]}{\end{enumerate}}

% Compatibility helpers needed by a few files in the combined archive.
\AtBeginDocument{\noindent{\footnotesize English translation prepared with ChatGPT (OpenAI, GPT-6), 26 October 2026, for translation of the Bulgarian text into English and LaTeX formatting.
The Bulgarian originals remain the reference; translation may contain errors.}\par\smallskip}
\newcommand{\MissingGraphic}[1]{%
  \begin{center}\fbox{\parbox{0.82\linewidth}{\centering Missing source image: \texttt{#1}}}\end{center}}



\pagestyle{fancy}\fancyhf{}\lhead{Applied Mathematics: Test 1 solutions}\rhead{20 April 2024}
\setlength{\headheight}{15pt}\pagenumbering{arabic}\cfoot{\thepage}
\begin{document}
\noindent\emph{English transcription of the original handwritten solution notes. The original calculations and arguments are retained. Apparent inconsistencies in the original are marked with numbered transcription notes at the end.}

\section*{Problem 1}
\subsection*{1. Conditional probability}
A yellow ball means rolling one die, green means two dice, and blue means three dice. Thus
\[
\P(\text{green}\mid\text{sum }3)
=\frac{\frac2{36}\cdot\frac13}
{\frac16\cdot\frac16+\frac2{36}\cdot\frac13+\frac1{216}\cdot\frac12}
=\frac8{21}\approx38.10\%.
\]
\subsection*{2. Five dice}
\[
\P(\max\ge3)=1-\P(\max=1\text{ or }2)=1-\left(\frac13\right)^5.
\]
For the probability that the product is divisible by 10, let $A$ be the event that at least one number is divisible by 5 and $B$ the event that at least one number is divisible by 2. Then
\[
\P(A^c)=\left(\frac56\right)^5,\qquad
\P(B^c)=\left(\frac12\right)^5,\qquad
\P(A^c\cap B^c)=\left(\frac26\right)^5.
\]
\begin{align*}
\P(A\cap B)&=\P(A)+\P(B)-\P(A\cup B)\\
&=1-\P(A^c)+1-\P(B^c)-\bigl(1-\P(A^c\cap B^c)\bigr)\\
&=1-\left(\frac56\right)^5-\left(\frac12\right)^5
+\left(\frac16\right)^5\approx56.70\%.\tag{Note 1}
\end{align*}

\subsection*{3. Seven balls in seven boxes}
For $X_3=2$, two boxes contain three balls each and one box contains one ball:
\[
\P(X_3=2)=\frac{7\binom73\,6\binom43\,5/2}{7^7}\approx1.78\%.
\]
For $X_3=1$, the possible nonzero occupancies are
$(3,4)$, $(3,1,1,1,1)$, $(3,2,2)$ and $(3,2,1,1)$. Hence
\begin{align*}
\P(X_3=1)&=\frac{7\binom73\,6}{7^7}
+\frac{7\binom73\,6\cdot5\cdot4\cdot3}{7^7}\\
&\quad+\frac{7\binom73\,6\binom42\,5/2}{7^7}
+\frac{7\binom73\,6\binom42\,5\cdot4}{7^7}
\approx34.99\%.
\end{align*}
The original then writes
\[
\P(X_3=0)=1-\P(X_3=0\text{ or }1)\approx63.23\%.\tag{Note 2}
\]
\paragraph{Check.}
Enumerate the arrangements with no box containing three balls. The following table transcribes the separate counts and their approximate contributions from the handwritten check. Every numerator is divided by $7^7$.
\begin{center}\small
\begin{tabular}{lll}
Nonzero occupancies & Numerator & Approximation\\\hline
$1+1+1+1+1+1+1$ & $7!$ & $0.61\%$\\[3pt]
$2+1+1+1+1+1$ & $7\binom72\,6\cdot5\cdot4\cdot3\cdot2$ & $12.85\%$\\[3pt]
$4+1+1+1$ & $7\binom74\,6\cdot5\cdot4$ & $3.57\%$\\[3pt]
$2+2+1+1+1$ & $7\binom72\,6\binom52\,5\cdot4\cdot3/2$ & $32.13\%$\\[3pt]
$5+1+1$ & $7\binom75\,6\cdot5$ & $0.33\%$\\[3pt]
$4+2+1$ & $7\binom74\,6\binom32\,5$ & $2.68\%$\\[3pt]
$2+2+2+1$ & $\binom73\frac{7\cdot6\cdot5\cdot4\cdot3\cdot2}{(2!)^3}\,4$ & $10.71\%$\\[3pt]
$6+1$ & $7\binom76\,6$ & $0\%$\\[3pt]
$5+2$ & $7\binom75\,6$ & $0\%$\\\hline
Total & & $63.23\%$
\end{tabular}
\end{center}
\[
\E X_3=\P(X_3=1)+2\P(X_3=2)
\approx\frac{35}{100}+\frac4{100}\approx\frac{37}{100}.\tag{Note 3}
\]
Also, in the notation of the original,
\[
\sum_{i=0}^7i\E X_i=\E[X_0+1X_1+\cdots+7X_7]=7,\tag{Note 4}
\]
because the weighted occupancy counts give the total number of balls.
Similarly,
\[
\sum_{i=0}^7\E X_i=\E[X_0+\cdots+X_7]
=\E[\text{number of boxes}]=7.
\]

\subsection*{4. Products and quotients}
For independent random variables, $\E[XY]=\E X\,\E Y$. See the lectures for the proof.
In general, $\E[X/Y]\ne\E X/\E Y$. A counterexample can be constructed from simple random variables, for instance independent $X,Y$ with
\[
\P(X=1)=\P(X=2)=\tfrac12,\qquad
\P(Y=2)=\P(Y=3)=\tfrac12.
\]

\section*{Problem 2}
\subsection*{1. A randomly labelled die}
Let $X$ be the number of faces labelled 6. Then $X\sim\operatorname{Bin}(6,1/6)$.
Let $Y$ be the outcome of the first roll. By total probability,
\begin{align*}
\P(Y=6)&=\sum_{k=0}^6\P(Y=6\mid X=k)\P(X=k)\\
&=\sum_{k=0}^6\frac k6\binom6k\left(\frac16\right)^k
\left(\frac56\right)^{6-k}
=\frac16\E[\operatorname{Bin}(6,1/6)]=\frac16.
\end{align*}
Let $Z$ be the number of rolls until the first 6. A die with no face labelled 6 is possible. Conditional on $X=0$, we have $Z=\infty$: a 6 is never rolled. Consequently, $\E Z=\infty$.

\subsection*{2. A product divisible by 6}
The expectation is again infinite, because all the faces might be labelled 5.

\subsection*{3. Two observed sixes}
Bayes' formula gives
\begin{align*}
&\P(\text{all faces are 6}\mid\text{two observed 6s})\\
&\qquad=\frac{1\cdot1\cdot\P(X=6)}
{\sum_{k=0}^6\P(X=k)\P(\text{two observed 6s}\mid X=k)}\\
&\qquad=\frac{(1/6)^6}
{\sum_{k=0}^6\binom6k(1/6)^k(5/6)^{6-k}(k/6)^2}\\
&\qquad=\frac{(1/6)^6}{(1/6)^2\E[\operatorname{Bin}(6,1/6)^2]}\\
&\qquad=\frac{(1/6)^4}{6\cdot(1/6)\cdot(5/6)+(6/6)^2}
\approx0.04\%.
\end{align*}
For the most likely number of sixes, compare
\[
f(k)=\P(X=k)(k/6)^2\ \propto\ \binom6k5^{6-k}(k/6)^2.
\]
The respective rescaled values are
\[
\begin{array}{c|rrrrrrr}
k&0&1&2&3&4&5&6\\\hline
\text{rescaled }f(k)&0&520.83&1041.67&625&166.67&20.83&1
\end{array}
\]
Thus, two faces labelled 6 is the most likely case.

\section*{Problem 3}
\subsection*{1. The probability of selecting $k$}
Write $p_n(k)=\P(k\in E_n)$, the probability of selecting $k$ when the first choice is from $\{1,\ldots,n\}$. The key observation is
\[
p_n(k)=p_n(k+1)+p_n(k+1)\frac1k.
\]
The first term corresponds to selecting $k$ without selecting $k+1$ (why?); the second corresponds to selecting $k+1$ and then $k$. Hence
\[
p_n(k)=\frac{k+1}{k}p_n(k+1).
\]
Using $p_n(n)=1/n$ (or $p_n(1)=1$), telescoping gives
\[
p_n(k)=\frac1k.
\]

\subsection*{2. Conditional probability}
\[
\P(2\in E_n\mid3\notin E_n)
=\frac{p_n(3)}{1-p_n(3)}=\frac{1/3}{2/3}=50\%.
\]

\subsection*{3. Mean and variance}
Let $X_i=\mathbf1_{\{i\in E_n\}}$. Then
\[
\E|E_n|=\E\sum_{i=1}^nX_i
=\sum_{i=1}^n\P(i\in E_n)=1+\frac12+\cdots+\frac1n\sim\ln n.
\]
For the variance, use
\[
D\left(\sum_{i=1}^nX_i\right)=\sum_{i,j=1}^n\operatorname{Cov}(X_i,X_j),
\qquad DX_i=\frac1i-\frac1{i^2}.
\]
For $i<j$,
\[
\operatorname{Cov}(X_i,X_j)=\E[X_iX_j]-\E X_i\E X_j
=\P(i,j\in E_n)-\frac1i\frac1j.
\]
However,
\[
\P(i,j\in E_n)=\P(i\in E_n\mid j\in E_n)\P(j\in E_n)
=\frac1i\frac1j\qquad\text{(why?)}.
\]
Thus the covariances vanish for $i\ne j$, and
\[
D|E_n|=\sum_{i=1}^nDX_i
=1+\frac12+\cdots+\frac1n
-\left(1+\frac1{2^2}+\cdots+\frac1{n^2}\right)\sim\ln n,
\]
since the sum of $1/i^2$ converges.

\section*{Problem 4}
\subsection*{1. Joint distribution and correlation}
The joint distribution is
\[
\begin{array}{c|ccc|c}
Z\backslash X&0&1&2&\P(Z=z)\\\hline
0&1/9&0&2/9&1/3\\
1&2/9&1/9&0&1/3\\
2&0&2/9&1/9&1/3\\\hline
\P(X=x)&1/3&1/3&1/3&1
\end{array}
\]
Therefore $\E X=\E Z=1$ and
\[
DX=DZ=\tfrac13\cdot1^2+\tfrac13\cdot2^2-1=\tfrac23.
\]
\[
\operatorname{Cor}(X,Z)=\frac{\E[XZ]-\E X\E Z}{\sqrt{DX\,DZ}}
=\frac{1\cdot1\cdot\frac19+2\cdot1\cdot\frac29+2^2\cdot\frac19-1^2}
{\sqrt{4/9}}=0.
\]

\subsection*{2. Generating function of a die}
\[
g_X(s)=\tfrac16(s+\cdots+s^6),\qquad
\E X=g'_X(1)=1,\tag{Note 5}
\]
\[
DX=g''_X(1)+g'_X(1)-[g'_X(1)]^2.
\]

\subsection*{3. Generating functions and independence}
If $X$ and $Y$ are independent, $\P(X=k,Y=m)=\P(X=k)\P(Y=m)$, so
\begin{align*}
\E s^{X+Y}
&=\sum_{k=0}^\infty s^k\P(X+Y=k)\\
&=\sum_{k=0}^\infty s^k\sum_{m=0}^k\P(X=m,Y=k-m)\\
&=\sum_{k=0}^\infty s^k\sum_{m=0}^k\P(X=m)\P(Y=k-m)\\
&=\sum_{k=0}^\infty\sum_{m=0}^k s^m\P(X=m)s^{k-m}\P(Y=k-m)\\
&=\sum_{m=0}^\infty s^m\P(X=m)\sum_{k=m}^\infty s^{k-m}\P(Y=k-m)\\
&=g_X(s)g_Y(s).
\end{align*}
The order of summation is changed over the region $0\le m\le k$.

The converse is false! The distribution in part 1 is a counterexample. The original expands the product as
\[
g_X(s)g_Z(s)=(1+s+s^2)^2(1/3)^2
=\frac{s^4+2s^3+3s^2+2s+s}{9}.\tag{Note 6}
\]
Meanwhile,
\[
\begin{array}{c|ccccc}
X+Z&0&1&2&3&4\\\hline
\text{probability}&1/9&2/9&3/9&2/9&1/9
\end{array}
\]
Thus $g_{X+Z}(s)=g_X(s)g_Z(s)$, although $X,Z$ are dependent:
\[
\P(X=0,Z=2)=0\ne\P(X=0)\P(Z=2).
\]

\subsection*{4. Alternative dice}
Yes, it is possible. For a standard die,
\[
g_X(s)=\tfrac16(s+\cdots+s^6),\qquad
g_{X+Y}(s)=\tfrac1{36}(s+\cdots+s^6)^2.
\]
We want to express this as the product of two generating functions other than the standard one. Equivalently, factor
\[
(s+\cdots+s^6)^2=P(s)Q(s),\qquad P(1)=Q(1)=6.
\]
The initial list of conditions in the original has $P(0)=Q(0)=1$ (why?), whereas the following page uses $P(0)=Q(0)=0$. See Note 7.
\begin{align*}
(s+\cdots+s^6)^2
&=[s(s^2+s+1)(s^3+1)]^2\\
&=[s(s+1)(s^2+s+1)(s^2-s+1)]^2.
\end{align*}
At $s=1$ these four factors have values $1,2,3,1$, respectively. The conditions $P(0)=Q(0)=0$ and $P(1)=Q(1)=6$ require the factors $s(s+1)(s^2+s+1)$ in both polynomials.

\paragraph{Case 1.}
\[
P(s)=s(s+1)(s^2+s+1),\qquad Q(s)=P(s)(s^2-s+1)^2.
\]
Expanding,
\[
P(s)=s^4+2s^3+2s^2+s,\qquad Q(s)=s^8+s^6+s^5+s^4+s^3+s.
\]
The dice therefore have faces
\[
(1,2,2,3,3,4)\qquad\text{and}\qquad(1,3,4,5,6,8).
\]
\paragraph{Case 2.}
\[
P(s)=Q(s)=s(s+1)(s^2+s+1)(s^2-s+1),
\]
which gives two standard dice.

\section*{Transcription notes}
These notes identify apparent slips in the Bulgarian handwritten source; they are not part of its original solution text.
\begin{enumerate}
\item The last term on the final line is written $(1/6)^5$ in the source, but the preceding line gives $(2/6)^5$. The displayed approximation $56.70\%$ agrees with $(2/6)^5$.
\item The source writes ``$0$ or $1$''; the complement of $X_3=0$ is $X_3=1$ or $X_3=2$.
\item The source's approximate arithmetic $35/100+4/100\approx37/100$ is inconsistent. The exact displayed probabilities give $\E X_3=\P(X_3=1)+2\P(X_3=2)$.
\item The source includes $X_0$ in the weighted sum. Its coefficient should be zero when counting balls. The unweighted sum, which counts boxes, does include $X_0$.
\item The source writes $g'_X(1)=1$. For the generating function on that same line, the derivative at 1 is $7/2$.
\item The last term of the expanded numerator is written $s$ in the source. Expansion of $(1+s+s^2)^2$ gives constant term 1, in agreement with the subsequent distribution table.
\item The source initially writes $P(0)=Q(0)=1$, then uses $P(0)=Q(0)=0$. The latter is consistent with positive integer die labels and with the factors displayed.
\end{enumerate}
\end{document}
